\documentclass{situationist-preprint}

\title{Attunement as Disclosure}
\msid{PHE-21-a}
\msfamily{PHE-21}
\msdiscipline{Phenomenology}
\mstype{Principal manuscript}
\mskeywords{moods; attunement; Befindlichkeit; Heidegger; Ratcliffe; existential feelings; affective intentionality; knowledge}

\begin{document}
\maketitle

\begin{abstract}
Are moods ways of knowing? The dominant view in the philosophy of emotion treats moods as non-cognitive states that colour experience without informing it, or as distorting influences on perception and judgment. We argue that moods are ways of knowing: they are pre-reflective disclosures of one's situation, fallible but not arbitrary. The argument develops Heidegger's analysis of attunement in \S 29 of \emph{Being and Time} and its recent elaborations by Ratcliffe and Slaby. Moods present one's situation as a whole as being some way, so they have correctness conditions; they do so without being objects of reflection; they are responsive to features of the situation that reflection may fail to register, so they are not arbitrary; and they can present the situation as it is not, as in depression, so they are fallible. A state with these features, when it is correct and properly based, can ground knowledge of one's situation. The result is consistent with the analyses of profound boredom and of anxiety as intentional disclosures.
\end{abstract}

\section{The problem}

Emotions are now widely held to have intentional objects and correctness conditions: fear presents its object as dangerous and is fitting when it is \citep{Goldie2000, DeonnaTeroni2012, Tappolet2016}. Moods are a harder case. They lack particular objects, they persist without being about anything in particular, and they often arise from causes, such as fatigue or illness, that have nothing to do with how things are. A common conclusion is that moods are not ways of knowing at all but background states that bias perception and judgment, and at best indicate something about the subject's body.

Heidegger defended the opposite view. Attunement (Befindlichkeit), the existential of which moods are the ontic manifestation, discloses Dasein's thrownness, discloses being-in-the-world as a whole, and makes it possible for things to matter, so that anything can be encountered as threatening, inviting or indifferent \citep[\S 29]{Heidegger1927}. Moods are not inner states projected onto a neutral world; they are the way in which one finds oneself in a situation. Heidegger adds that moods disclose more primordially than cognition, and that the disclosure is not to be confused with knowing in the sense of theoretical cognition.

\citet{Ratcliffe2008} has developed this into an account of \emph{existential feelings}: bodily feelings that constitute a sense of reality and of belonging to the world, and that structure the space of possibilities in which particular experiences occur. \citet{Slaby2008} has argued for an affective intentionality in which bodily feeling is itself world-directed. On both accounts moods disclose something. Whether what they disclose is knowledge, and whether moods can be wrong, remain open. We argue that moods are ways of knowing in a precise sense: pre-reflective disclosures of one's situation, fallible but not arbitrary.

\section{Framework}

\begin{definition}[Way of knowing]\label{d:wok}
A kind of state is a \emph{way of knowing} if its instances present how things are, can be correct or incorrect, and, when correct and formed in a way that is responsive to what they present, can constitute or ground knowledge of it.
\end{definition}

The definition is deliberately broad: perception and testimony are ways of knowing in this sense, though both are fallible. It does not require that the state be a belief.

\begin{definition}[Pre-reflective disclosure]\label{d:pre}
A state is a \emph{pre-reflective disclosure} of a situation if it presents the situation without the situation, or the state itself, being an object of reflection or of explicit attention.
\end{definition}

\begin{principle}[Situational presentation]\label{p:pres}
A mood presents the subject's situation as a whole as being some way: as hopeful or closed, safe or threatening, inviting or flat.
\end{principle}

\begin{principle}[Responsiveness]\label{p:resp}
Moods, in their normal operation, vary with features of the subject's situation that bear on how things matter to her, including features that she has not explicitly registered.
\end{principle}

\begin{principle}[Error]\label{p:err}
A mood can present the situation as being a way that it is not.
\end{principle}

Principle~\ref{p:pres} is Heidegger's claim that attunement discloses being-in-the-world as a whole, stated without his vocabulary. PHE-04 argues that undirected moods are intentional in this sense, presenting the world or one's situation as a whole, and that even profound boredom is directed, towards time itself. Principle~\ref{p:resp} is supported by ordinary experience and by clinical observation: a vague unease in a meeting can register hostility before anyone could say what signalled it, and a lifting of mood can register, before reflection catches up, that a danger has passed. Principle~\ref{p:err} is supported by depression, in which the world is presented as devoid of possibility although it is not \citep{Ratcliffe2015}, and by moods induced by drugs or fatigue.

\section{The argument}

\begin{lemma}\label{l:corr}
Moods have correctness conditions.
\end{lemma}

\begin{proof}
By Principle~\ref{p:pres}, a mood presents the situation as being some way. A presentation of something as being some way is correct if it is that way and incorrect otherwise. By Principle~\ref{p:err}, the incorrect case occurs, so the condition is not empty.
\end{proof}

\begin{lemma}\label{l:pre}
Moods are pre-reflective.
\end{lemma}

\begin{proof}
A mood presents the situation by colouring how everything in it appears, not by making the situation an object of attention: in gloom, things look flat, and one need not attend to the gloom, or to the situation as such, for them to look so. The mood can become an object of reflection, as when one notices that one is irritable, but it was operative before, and noticing it does not create it. By Definition~\ref{d:pre}, moods are pre-reflective disclosures.
\end{proof}

\begin{lemma}\label{l:arb}
Moods are not arbitrary.
\end{lemma}

\begin{proof}
A kind of presentation is arbitrary if whether it is correct is unrelated to whether what it presents obtains. By Principle~\ref{p:resp}, moods in their normal operation vary with the features of the situation that bear on how things matter. Hence their correctness is related to how the situation is: when the situation becomes threatening, unease tends to arise, and when it ceases to be, unease tends to subside. Moods are therefore not arbitrary.
\end{proof}

\begin{theorem}\label{t:main}
Moods are ways of knowing: they are pre-reflective disclosures of one's situation, fallible but not arbitrary.
\end{theorem}

\begin{proof}
By Principle~\ref{p:pres} and Lemma~\ref{l:corr}, moods present how the situation is and can be correct or incorrect. By Lemma~\ref{l:pre}, they do so pre-reflectively. By Principle~\ref{p:err}, they are fallible; by Lemma~\ref{l:arb}, they are not arbitrary. It remains to show that a correct mood, formed responsively, can ground knowledge. Consider the unease that registers hostility in a meeting. If the hostility is there, and the unease arose because of the cues that manifest it, the subject who trusts her unease and concludes that the meeting is hostile has a belief that is true and is based on a state that is sensitive to its truth. That is as good a basis for knowledge as a perceptual experience formed in the same way. By Definition~\ref{d:wok}, moods are a way of knowing.
\end{proof}

\begin{remark}\label{r:object}
What moods disclose is how the situation as a whole matters, not facts about particular objects. They are therefore a way of knowing distinct from perception, and their characteristic objects include features that no particular perception presents: that one's possibilities are open or closed, that one's surroundings are or are not a place where one belongs. PHE-16 argues that anxiety discloses the withdrawal of beings as a whole, a genuine disclosure that reveals no further entity. Anxiety is the limit case of Theorem~\ref{t:main}, in which what is disclosed is that the significance of beings is not secured by them.
\end{remark}

\section{Objections and replies}

\begin{objection}[Causes, not reasons]
Moods are caused by hormones, fatigue, weather and illness. A state with such causes cannot be a way of knowing, since its occurrence is explained by facts unrelated to what it supposedly presents.
\end{objection}

\begin{reply}
Every way of knowing has non-cognitive causes; perception is caused by retinal stimulation. The question is whether the state's occurrence is, in normal operation, responsive to what it presents, and Principle~\ref{p:resp} says it is. A mood caused by fatigue alone is an instance of Principle~\ref{p:err}: fatigue presents the world as heavier than it is. That moods can be produced by causes that bypass the situation shows that they are fallible, which the theorem asserts, not that they are arbitrary, which it denies.
\end{reply}

\begin{objection}[Depression]
Depression shows that moods can present the world as wholly devoid of possibility, persistently and in the face of evidence. A way of knowing should be corrigible by evidence; a mood that resists all evidence is arbitrary.
\end{objection}

\begin{reply}
Depression is a disorder of the way of knowing, as hallucination is a disorder of perception. \citet{Ratcliffe2015} describes it as an alteration of the space of possibilities within which experience is structured, so that the depressed person cannot experience things as offering hope. The persistence of the mood in the face of evidence is part of the disorder, not evidence that moods in general are unrelated to their objects. Moreover, Ratcliffe's description presupposes Principle~\ref{p:pres}: the depressed person's world is presented as closed, and it is because it is not closed that the presentation is in error.
\end{reply}

\begin{objection}[Hijacked perception]
\citet{Siegel2017} argues that moods and other states can influence perceptual experience in ways that make it epistemically worse, as when fear makes a neutral face look angry. Moods are sources of distortion of other ways of knowing, not ways of knowing themselves.
\end{objection}

\begin{reply}
The objection concerns the effect of a mood on the perception of particular objects. Theorem~\ref{t:main} concerns what moods present of the situation as a whole (Remark~\ref{r:object}). A fear that makes a neutral face look angry is incorrect as a presentation of the face, and it may also be incorrect as a presentation of the situation, if the situation is not threatening; then it is a case of Principle~\ref{p:err}. Where the situation is threatening, the fear may be correct about it while distorting the perception of a particular face. That a source of knowledge can interfere with another is familiar from the interaction of the senses.
\end{reply}

\begin{objection}[Non-conceptual content]
Knowledge is propositional. Moods do not have propositional content, so they cannot be knowledge, and they can ground knowledge only by being taken up into judgment, at which point it is the judgment that knows.
\end{objection}

\begin{reply}
Definition~\ref{d:wok} requires only that a correct and responsively formed mood can constitute or ground knowledge. Perceptual experiences ground knowledge in the same way on many views, and no one concludes that perception is not a way of knowing. Heidegger's own point that attunement is not cognition in the theoretical sense \citep[\S 29]{Heidegger1927} is consistent with this: the disclosure precedes and makes possible explicit cognition of one's situation, and it is for that reason that it can ground it.
\end{reply}

\section{Conclusion}

Moods present one's situation as a whole as being some way, without being objects of reflection. They vary, in normal operation, with what they present, and they can present it wrongly. They are therefore ways of knowing in the sense in which perception is: pre-reflective disclosures of one's situation, fallible but not arbitrary.

\begin{thebibliography}{99}\small
\bibitem[Deonna and Teroni(2012)]{DeonnaTeroni2012} Deonna, J.~A. and Teroni, F. (2012). \emph{The Emotions: A Philosophical Introduction}. Routledge.
\bibitem[Goldie(2000)]{Goldie2000} Goldie, P. (2000). \emph{The Emotions: A Philosophical Exploration}. Oxford University Press.
\bibitem[Heidegger(1927)]{Heidegger1927} Heidegger, M. (1927). \emph{Sein und Zeit}. Niemeyer. Trans. J. Macquarrie and E. Robinson, \emph{Being and Time}, Harper \& Row, 1962.
\bibitem[Ratcliffe(2008)]{Ratcliffe2008} Ratcliffe, M. (2008). \emph{Feelings of Being: Phenomenology, Psychiatry and the Sense of Reality}. Oxford University Press.
\bibitem[Ratcliffe(2015)]{Ratcliffe2015} Ratcliffe, M. (2015). \emph{Experiences of Depression: A Study in Phenomenology}. Oxford University Press.
\bibitem[Siegel(2017)]{Siegel2017} Siegel, S. (2017). \emph{The Rationality of Perception}. Oxford University Press.
\bibitem[Slaby(2008)]{Slaby2008} Slaby, J. (2008). Affective intentionality and the feeling body. \emph{Phenomenology and the Cognitive Sciences}, 7(4), 429--444.
\bibitem[Tappolet(2016)]{Tappolet2016} Tappolet, C. (2016). \emph{Emotions, Values, and Agency}. Oxford University Press.
\end{thebibliography}

\end{document}
